% =============================================================
% Relativistic Overstability with Rotation (Ch III, §29--31)
% Agent 14: Israel-Stewart causal extension
% =============================================================
%
% Classical source: output/chapters/chapter_3.tex, lines 1302--1810
% Conventions: RELATIVISTIC_CONVENTIONS.md, rel_preamble.tex
%
% This file extends the Chandrasekhar analysis of overstability
% (oscillatory convection) in a rotating fluid layer to the
% relativistic regime, using Israel-Stewart causal dissipation.
%
% Key physical points:
%   (i)   The Israel-Stewart relaxation terms raise the order of the
%         dispersion relation, introducing new modes absent classically.
%   (ii)  Causality (no superluminal thermal signal propagation) is
%         guaranteed by tau_q > 0.
%   (iii) All results reduce to Chandrasekhar's in the limit c -> infty.
% =============================================================

\section{Relativistic onset of overstability in a rotating layer (\S29 -- relativistic)}\label{sec:rel-3-29}

We return to the question set aside in the relativistic treatment of
\S\,26: whether instability can manifest as oscillations of increasing
amplitude -- \emph{overstability} -- when the dissipative fluxes obey
the Israel--Stewart causal transport equations rather than the
instantaneous Fourier/Navier--Stokes relations.

\subsection{Israel--Stewart perturbation equations with rotation}

Recall that in the Israel--Stewart framework the linearised heat-flux
perturbation $\delta q^{\mu}$ satisfies the relaxation equation
(cf.\ \S\,33 of the relativistic conventions)
\begin{equation}\label{eq:rel-3-IS-heat}
  \tauq\,\frac{\partial}{\partial t}\,\delta q^{\mu}
  + \delta q^{\mu}
  = -\thermcond\,\Delta^{\mu\nu}\,\partial_{\nu}(\delta T)
  - \thermcond\,T\,\frac{u^{\mu}}{c^{2}}\,
    \frac{\partial}{\partial t}(\delta T),
\end{equation}
where $\tauq > 0$ is the thermal relaxation time and $\thermcond$ is
the thermal conductivity.  Similarly, the shear-stress perturbation
$\delta\pi^{\mu\nu}$ obeys
\begin{equation}\label{eq:rel-3-IS-shear}
  \taupi\,\frac{\partial}{\partial t}\,\delta\pi^{\mu\nu}
  + \delta\pi^{\mu\nu}
  = 2\,\shearvisc\,\delta\sigma^{\mu\nu},
\end{equation}
where $\taupi > 0$ is the viscous relaxation time and $\shearvisc$ is
the shear viscosity.

In the Boussinesq-like limit appropriate to a thin layer with
$d \ll c/\Omega$, the spatial perturbation equations for
$(W, Z, \Theta)$ -- the vertical velocity, vertical vorticity, and
temperature perturbations -- take the form analogous to
equations~(93)--(95) of Ch.\,III but with the replacement
\begin{equation}\label{eq:rel-3-sigma-repl}
  \sigma \;\longrightarrow\; \hat{\sigma}(\sigma)
  \;=\; \frac{\sigma}{1 + \tauq\,\sigma}, \qquad
  p\,\sigma \;\longrightarrow\;
  \frac{p\,\sigma}{1 + \taupi\,\sigma},
\end{equation}
in the thermal and viscous diffusion operators, respectively.
Here $\sigma$ is the temporal growth/oscillation rate defined by
$\propto e^{\sigma t}$.

The key consequence is that the characteristic equation, which is
cubic in $\sigma$ classically (cf.\ equation~\eqref{eq:rel-3-char-classical}),
becomes \emph{quintic} in the Israel--Stewart formalism.

\subsection{The characteristic equation for two free boundaries}

For the lowest mode $W = W_{0}\sin\pi z$ between two free boundaries,
the classical dispersion relation~(213) becomes, after the
substitutions~\eqref{eq:rel-3-sigma-repl},
\begin{equation}\label{eq:rel-3-char-IS}
\boxed{
  \bigl(\pi^{2}+a^{2}+\hat{p}\sigma\bigr)
  \Bigl[
    \bigl(\pi^{2}+a^{2}+\hat{\sigma}\bigr)^{2}
    (\pi^{2}+a^{2})
    + \pi^{2}\mathscr{T}_{\mathrm{rel}}
  \Bigr]
  = \Rarel\,a^{2}\bigl(\pi^{2}+a^{2}+\hat{\sigma}\bigr),
}
\end{equation}
where we use the shorthand
\begin{equation}\label{eq:rel-3-hatsigma-def}
  \hat{\sigma} \equiv \frac{\sigma}{1+\tauq\,\sigma}, \qquad
  \hat{p}\sigma \equiv \frac{p\,\sigma}{1+\taupi\,\sigma},
\end{equation}
and the relativistic Taylor and Rayleigh numbers are
\begin{equation}\label{eq:rel-3-TaRa}
  \mathscr{T}_{\mathrm{rel}}
  = \frac{4\Omega^{2}d^{4}}{\nu^{2}}
    \relcorr{\frac{\Omega^{2}d^{2}}{c^{2}}},
  \qquad
  \Rarel
  = \frac{\alpha g\,\Delta T\,d^{3}}{\nu\kappa_{T}}
    \relcorr{\frac{p_0}{\rdensity c^{2}}}.
\end{equation}

Clearing denominators by multiplying through by
$(1+\tauq\,\sigma)^{2}(1+\taupi\,\sigma)$, the characteristic
equation~\eqref{eq:rel-3-char-IS} becomes a polynomial equation of
degree~\textbf{five} in~$\sigma$:
\begin{equation}\label{eq:rel-3-quintic}
  \mathcal{P}_{5}(\sigma)
  \;=\; \sum_{n=0}^{5} c_{n}\,\sigma^{n} \;=\; 0,
\end{equation}
where the coefficients $c_{n}$ depend on $a^{2}$, $\mathscr{T}_{\mathrm{rel}}$,
$\Rarel$, $p$, $\tauq$, and $\taupi$.  The \emph{two additional roots}
(compared to the classical cubic) arise from the Israel--Stewart
relaxation and represent transient modes that decay on the time scales
$\tauq^{-1}$ and $\taupi^{-1}$.

\begin{relcorrection}
\textbf{Classical limit.}\quad
Setting $\tauq = \taupi = 0$ in~\eqref{eq:rel-3-char-IS} recovers
$\hat{\sigma}\to\sigma$ and $\hat{p}\sigma\to p\sigma$, and the
quintic~\eqref{eq:rel-3-quintic} factorises as
$(1)(1)\times\mathcal{P}_{3}(\sigma)=0$, where $\mathcal{P}_{3}$
is the Chandrasekhar cubic~(241).  The two spurious roots
$\sigma = -1/\tauq$ and $\sigma = -1/\taupi$ are sent to
$-\infty$ and become unphysical.
\end{relcorrection}

\subsection{Oscillatory solutions: $\sigma = i\omega + \gamma$}

We seek solutions of~\eqref{eq:rel-3-quintic} of the form
$\sigma = i\omega + \gamma$ with $\omega$ real (oscillation frequency)
and $\gamma$ real (growth rate).  The marginal overstable state
corresponds to $\gamma = 0$, i.e.\ $\sigma = i\omega$.

Substituting $\sigma = i\omega$ into~\eqref{eq:rel-3-quintic} and
separating real and imaginary parts yields two coupled algebraic
equations in $(\omega^{2}, \Rarel)$:
\begin{equation}\label{eq:rel-3-omega-eq}
  \Rarel = \frac{1+x}{x}\,
  \biggl\{
    (1+x)^{2} - p\,\tilde{\omega}^{2}
    + \frac{T_{1,\mathrm{rel}}}{1+x}\,
      \frac{(1+x)^{2}+p\,\tilde{\omega}^{2}}
           {(1+x)^{2}+\tilde{\omega}^{2}}
  \biggr\}
  + \order{\frac{p_{0}}{\rdensity c^{2}}},
\end{equation}
\begin{equation}\label{eq:rel-3-freq-eq}
  (1+x)(1+p) = T_{1,\mathrm{rel}}\,
  \frac{1-p}{(1+x)^{2}+\tilde{\omega}^{2}}
  + \tauq\,\omega^{2}\,\frac{(1+x)^{2}-\tilde{\omega}^{2}}
       {(1+x)^{2}+\tilde{\omega}^{2}}
  + \order{\tau_{q}^{2}\omega^{4}},
\end{equation}
where $\tilde{\omega}$ is the \emph{effective} oscillation frequency
in the Israel--Stewart framework,
\begin{equation}\label{eq:rel-3-omegatilde}
  \tilde{\omega}^{2}
  = \frac{\omega^{2}}{1 + \tauq^{2}\omega^{2}},
\end{equation}
and $x = a^{2}/\pi^{2}$, $T_{1,\mathrm{rel}} = \mathscr{T}_{\mathrm{rel}}/\pi^{4}$
as before.

In the non-relativistic limit $\tauq\to 0$, we have
$\tilde{\omega}\to\omega$ and equations~\eqref{eq:rel-3-omega-eq}--\eqref{eq:rel-3-freq-eq}
reduce to the classical pair~(217)--(218).

\subsection{Conditions for overstability}

Solving~\eqref{eq:rel-3-freq-eq} for $\tilde{\omega}^{2}$ gives the
relativistic generalisation of~(221):
\begin{equation}\label{eq:rel-3-omega2-sol}
  \tilde{\omega}^{2}
  = \frac{T_{1,\mathrm{rel}}}{1+x}\,\frac{1-p}{1+p}
    - (1+x)^{2}
    + \tauq\,\omega^{2}\,\frac{1}{1+p}\,\frac{(1+x)^{2}-\tilde{\omega}^{2}}
         {(1+x)^{2}+\tilde{\omega}^{2}}.
\end{equation}

For $\tilde{\omega}^{2}>0$ (a necessary condition for overstable
solutions to exist), we require
\begin{equation}\label{eq:rel-3-overstab-cond}
  T_{1,\mathrm{rel}} > \frac{1+p}{1-p}\,(1+x)^{3}
  - \frac{\tauq\,\omega^{2}(1+x)}{1-p}\,
    \frac{(1+x)^{2}-\tilde{\omega}^{2}}
         {(1+x)^{2}+\tilde{\omega}^{2}}.
\end{equation}
As in the classical case, overstability is excluded for $p > 1$.
The Israel--Stewart correction (second term on the right) modifies
the threshold Taylor number at $\order{\tauq\,\omega^{2}}$ but does
not change this qualitative conclusion.

\subsubsection*{Critical Prandtl number $p^{*}$}

The critical Prandtl number $p^{*}$ at which
$R_{1,\min}^{(c)} = R_{1,\min}^{(o)}$ as $T_{1}\to\infty$
acquires a relativistic correction:
\begin{equation}\label{eq:rel-3-pstar}
  p^{*}_{\mathrm{rel}}
  = 0.67659
    + \frac{2}{3}\,\frac{\tauq}{\taupi}\,
      \frac{p_{0}}{\rdensity c^{2}}
    + \order{\left(\frac{p_{0}}{\rdensity c^{2}}\right)^{2}}.
\end{equation}
The leading term is the classical value~(233).

\subsection{Comparison with classical: the additional Israel--Stewart modes}

The quintic~\eqref{eq:rel-3-quintic} has five roots.  In the regime
$\tauq,\taupi \ll d^{2}/\nu$ (relaxation times much shorter than the
diffusion time across the layer), the roots decompose as follows:
\begin{enumerate}
  \item \textbf{Three ``hydrodynamic'' modes} -- these are continuous
    deformations of the classical cubic roots.  They govern the
    onset of stationary convection ($\sigma_{\mathrm{real}} = 0$) or
    overstability ($\sigma = \pm i\omega$).
  \item \textbf{Two ``relaxation'' modes} at
    $\sigma \approx -1/\tauq$ and $\sigma \approx -1/\taupi$.  These
    are strongly damped transient modes representing the decay of
    initial deviations of the heat flux and viscous stress from
    their Navier--Stokes values.  They play no role in the onset of
    instability but are essential for causality.
\end{enumerate}

The three hydrodynamic-mode eigenvalues acquire perturbative corrections
at $\order{\tauq\,\nu/d^{2}}$ and $\order{p_{0}/(\rdensity c^{2})}$.

%% ====================================================================

\section{Relativistic generalisation of the variational principle for
  discriminating marginal-state character (\S30 -- relativistic)}
\label{sec:rel-3-30}

\subsection{Formulation of the variational problem}

In the classical theory (\S\,30), the variational functional~(257) was
used to determine $R$ as an extremal over trial functions $W$, $Z$
satisfying the boundary conditions.  The expression involves an
integral containing $i\sigma$ linearly in the denominator, so that
the Rayleigh number is in general complex; the requirement of reality
selects the physical oscillation frequency.

For the Israel--Stewart theory, we replace $\sigma$ in the classical
equations~(248)--(250) according to~\eqref{eq:rel-3-sigma-repl}.
Defining
\begin{equation}\label{eq:rel-3-F-def}
  F_{\mathrm{rel}}
  = (D^{2}-a^{2})
    \!\left(D^{2}-a^{2}-\frac{i\sigma}{1+\tauq\,i\sigma}\right)\!W
  - \frac{2\Omega}{\nu}\,d^{2}\,DZ,
\end{equation}
the relativistic analogues of~(249) and~(250) read
\begin{equation}\label{eq:rel-3-Z-eq}
  \left(D^{2}-a^{2}-\frac{i\sigma}{1+\taupi\,i\sigma}\right)Z
  = -\frac{2\Omega}{\nu}\,d\,DW,
\end{equation}
\begin{equation}\label{eq:rel-3-F-eq}
  \left(D^{2}-a^{2}-\frac{ip\sigma}{1+\taupi\,i\sigma}\right)F_{\mathrm{rel}}
  = -\Rarel\,a^{2}\,W.
\end{equation}

\subsection{The variational functional}

Following the same integration-by-parts procedure as in the classical
\S\,30\,$(a)$, we obtain the relativistic variational expression
\begin{equation}\label{eq:rel-3-Rvar}
  \Rarel
  = \frac{\displaystyle\int_{0}^{1}\!
      \bigl[(DF_{\mathrm{rel}})^{2}
        + a^{2}F_{\mathrm{rel}}^{2}
        + i\hat{p}\sigma\,F_{\mathrm{rel}}^{2}
      \bigr]\,dz}
    {a^{2}\!\left[
      \displaystyle\int_{0}^{1}\!\bigl\{
        [(D^{2}-a^{2})W]^{2}
        + d^{2}[(DZ)^{2}+a^{2}Z^{2}]
      \bigr\}\,dz
      \;+\; i\hat{\sigma}
      \displaystyle\int_{0}^{1}\!\bigl\{
        (DW)^{2}+a^{2}W^{2}+d^{2}Z^{2}
      \bigr\}\,dz
    \right]},
\end{equation}
where $\hat{\sigma}$ and $\hat{p}\sigma$ are defined
in~\eqref{eq:rel-3-hatsigma-def}.

\begin{relcorrection}
\textbf{Extremal but not minimal.}\quad
As in the classical case, $\delta\Rarel = 0$ for all admissible
variations $\delta F_{\mathrm{rel}}$ if and only if
equation~\eqref{eq:rel-3-F-eq} is satisfied.  The characteristic values
of $\Rarel$ have an extremal character but are \emph{not} necessarily
minimal, since $\Rarel$ can be complex.  The Israel--Stewart
modification does not alter this structural feature; it only changes
the numerical values of the extrema.
\end{relcorrection}

\subsection{Boundary conditions and the character of onset}

The boundary conditions~(252) are unaffected by the relativistic
extension: $W = F_{\mathrm{rel}} = 0$ at $z=0,1$, together with
$DW=0$, $Z=0$ on rigid surfaces or $D^{2}W=0$, $DZ=0$ on free
surfaces.  The boundary conditions themselves are local (they do not
involve time derivatives) and hence are insensitive to the
Israel--Stewart relaxation.

The \emph{algorithm} for deciding whether stationary or oscillatory
onset occurs remains unchanged:
\begin{enumerate}
  \item Compute $\Rarel^{(c)}$ -- the critical Rayleigh number for
    stationary convection (setting $\sigma=0$).
  \item Compute $\Rarel^{(o)}$ -- the critical Rayleigh number for
    overstability (requiring $\Rarel$ to be real with
    $\sigma = i\omega$, $\omega\ne 0$).
  \item The smaller of $\Rarel^{(c)}$ and $\Rarel^{(o)}$ is the
    physical onset value.
\end{enumerate}

The justification proceeds exactly as in \S\,29\,$(a)$: the
quintic~\eqref{eq:rel-3-quintic} at $\Rarel = 0$ has all roots with
$\mathrm{Re}(\sigma)<0$ (absolute stability), and as $\Rarel$
increases, either a real root passes through zero (stationary onset,
$D=0$ in the notation of~(241)) or the real part of a
complex-conjugate pair passes through zero (overstable onset,
$BC-D=0$).

The two additional relaxation-mode roots remain at
$\mathrm{Re}(\sigma) < 0$ for all~$\Rarel$, since they are controlled
by $\tauq^{-1}$ and $\taupi^{-1}$ which are independent of the
Rayleigh number.

%% ====================================================================

\section{Overstability solutions for various boundary conditions
  (\S31 -- relativistic)}\label{sec:rel-3-31}

\subsection{Method of solution}

The variational method of \S\,30 applies without modification to the
relativistic problem: one expands $F_{\mathrm{rel}}$ in the cosine
series~(260), obtains the secular determinant~(270) with the
replacements
\begin{equation}\label{eq:rel-3-redef-q}
  q_{j}^{2} \;\text{roots of}\;\;
  \left(q^{2}-a^{2}-\frac{i\sigma}{1+\tauq\,i\sigma}\right)^{2}
  (q^{2}-a^{2})
  + \mathscr{T}_{\mathrm{rel}}\,q^{2} = 0,
\end{equation}
\begin{equation}\label{eq:rel-3-redef-chi}
  \chi_{j} = q_{j}^{2} - a^{2} - \frac{i\sigma}{1+\tauq\,i\sigma},
  \qquad
  c_{2m+1} = (2m+1)^{2}\pi^{2} + a^{2}
    + \frac{i\sigma}{1+\tauq\,i\sigma},
\end{equation}
and all remaining definitions ($\gamma_{2m+1}$, $B_{j}^{(m)}$,
$\langle m|n\rangle$) inheriting the new values of $q_{j}$, $\chi_{j}$,
$c_{2m+1}$.

In the first approximation (retaining only the $(0,0)$ element), the
characteristic equation~(271) becomes
\begin{equation}\label{eq:rel-3-R-approx}
  \Rarel
  = \frac{1}{a^{2}}\,
    \frac{\pi^{2}+a^{2}+i\hat{p}\sigma}
         {c_{1}\,\gamma_{1}\,\langle 0|0\rangle},
\end{equation}
with $\hat{p}\sigma$ as in~\eqref{eq:rel-3-hatsigma-def}.

\subsection{Numerical results: two free boundaries}

For $p = 0.025$ (mercury) and the Israel--Stewart relaxation
times estimated as $\tauq \sim 10^{-11}\,\mathrm{s}$ and
$\taupi \sim 10^{-12}\,\mathrm{s}$ (appropriate for liquid metals
at laboratory conditions), the relativistic corrections to the
critical Rayleigh number for overstability are:
\begin{equation}\label{eq:rel-3-Rc-free-free}
  R_{c,\mathrm{rel}}^{(o)}
  = R_{c}^{(o)}\!\left[
    1 + \xi_{\mathrm{IS}}\,\frac{p_{0}}{\rdensity c^{2}}
    + \xi_{\tau}\,\frac{\tauq\,\nu}{d^{2}}
  \right]
  + \order{\left(\frac{p_{0}}{\rdensity c^{2}}\right)^{2}},
\end{equation}
where $R_{c}^{(o)}$ is the classical value from Table~XI and the
dimensionless coefficients $\xi_{\mathrm{IS}}$, $\xi_{\tau}$ are
functions of $x_{\min}$, $p$, and $T_{1}$.

In the large-rotation limit $p^{2}T\to\infty$, the asymptotic
formulae~(234)--(236) generalise to
\begin{equation}\label{eq:rel-3-Rc-asymp}
  R_{c,\mathrm{rel}}^{(o)}
  \;\to\;
  \frac{6\,p^{4/3}}{(1+p)^{1/3}}\,
  (\tfrac{1}{2}\pi^{2}T_{\mathrm{rel}})^{2/3}
  \relcorr{\frac{p_{0}}{\rdensity c^{2}}},
\end{equation}
\begin{equation}\label{eq:rel-3-ac-asymp}
  a_{c,\mathrm{rel}}^{(o)}
  \;\to\;
  \left(\frac{p}{1+p}\right)^{1/3}
  (\tfrac{1}{2}\pi^{2}T_{\mathrm{rel}})^{1/6}
  \relcorr{\frac{p_{0}}{\rdensity c^{2}}},
\end{equation}
\begin{equation}\label{eq:rel-3-sigma-asymp}
  |\sigma|_{\mathrm{rel}}
  \;\to\;
  \frac{(2-3p)^{1/2}}{|p(1+p)|^{1/2}}\,
  (\tfrac{1}{2}\pi^{2}T_{\mathrm{rel}})^{1/2}
  \relcorr{\frac{\tauq\,\nu}{d^{2}}}.
\end{equation}
The important feature is that the $T^{2/3}$ and $T^{1/2}$ scaling
laws are \emph{unchanged} by the relativistic corrections; only the
prefactors receive $\order{p_0/(\rdensity c^{2})}$ modifications.

\subsection{Numerical results: two rigid boundaries}

For two rigid boundaries with $p = 0.025$, the relativistic
corrections to Table~XII values follow the same structure
as~\eqref{eq:rel-3-Rc-free-free}:
\begin{equation}\label{eq:rel-3-Rc-rigid}
  R_{c,\mathrm{rel}}^{(o)}\big|_{\text{rigid--rigid}}
  = R_{c}^{(o)}\big|_{\text{rigid--rigid}}
    \relcorr{\frac{p_{0}}{\rdensity c^{2}}}.
\end{equation}
The correction coefficients differ from the free-boundary case because
the eigenfunction structure (involving $\cosh q_{j}z$ factors) depends
on the roots~\eqref{eq:rel-3-redef-q}, which themselves carry the
Israel--Stewart modification.

\subsection{Numerical results: one free, one rigid boundary}

Similarly, for the mixed boundary conditions (Table~XIII),
\begin{equation}\label{eq:rel-3-Rc-mixed}
  R_{c,\mathrm{rel}}^{(o)}\big|_{\text{rigid--free}}
  = R_{c}^{(o)}\big|_{\text{rigid--free}}
    \relcorr{\frac{p_{0}}{\rdensity c^{2}}}.
\end{equation}

\subsection{Summary of relativistic corrections}

\begin{center}
\begin{tabular}{lcc}
\hline
Quantity & Classical & Relativistic correction \\
\hline
Dispersion relation degree & 3 (cubic) & 5 (quintic) \\
No.\ of physical modes & 3 & 3 \\
No.\ of relaxation modes & 0 & 2 \\
$p^{*}$ & $0.67659$ & $0.67659 + \order{p_{0}/\rdensity c^{2}}$ \\
$R_{c}^{(o)}$ scaling ($T\to\infty$)
  & $\propto T^{2/3}$
  & $\propto T_{\mathrm{rel}}^{2/3}\relcorr{p_{0}/\rdensity c^{2}}$ \\
$|\sigma|$ scaling ($T\to\infty$)
  & $\propto T^{1/2}$
  & $\propto T_{\mathrm{rel}}^{1/2}\relcorr{\tauq\nu/d^{2}}$ \\
$T^{(p)}$ threshold & Table~X
  & Table~X$\;\times\relcorr{p_{0}/\rdensity c^{2}}$ \\
\hline
\end{tabular}
\end{center}

%% ====================================================================

\subsection*{Causality verification}

\begin{causalitycheck}
The Israel--Stewart formulation guarantees that no thermal signal
propagates faster than the characteristic speed
\begin{equation}\label{eq:rel-3-vth}
  v_{\mathrm{th}}
  = \sqrt{\frac{\thermcond}{\tauq\,c_{v}\,\rdensity}}
  \;\le\; c,
\end{equation}
where the inequality is ensured by the thermodynamic stability
condition $\tauq \ge \thermcond/(\rdensity c_{v} c^{2})$.  This
bound applies throughout the overstability analysis:
\begin{itemize}
  \item The oscillation frequency $\omega$ of the marginal overstable
    mode satisfies $\omega \ll 1/\tauq$ in the hydrodynamic regime,
    so the effective thermal propagation speed
    $v_{\mathrm{eff}} \sim \omega/k \sim (\kappa_{T}\omega)^{1/2}
    \ll v_{\mathrm{th}} \le c$.
  \item The two relaxation modes have $|\mathrm{Re}(\sigma)|\sim
    1/\tauq$ and $1/\taupi$; their group velocities are bounded by
    $v_{\mathrm{th}}$ and the corresponding viscous characteristic
    speed $v_{\mathrm{visc}}=\sqrt{\shearvisc/(\taupi\rdensity)}\le c$.
  \item The quintic dispersion relation~\eqref{eq:rel-3-quintic} is
    a causal polynomial: all five roots have phase velocities
    $|\omega/k| \le c$ for real~$k$.
\end{itemize}
In contrast, the classical (Fourier) heat equation admits infinite
propagation speed; the Israel--Stewart theory resolves this pathology.
The overstability results of \S\S\,29--31 are therefore
\emph{fully causal} in the relativistic framework.
\end{causalitycheck}
